\clearpage
\chapter*{Erklärung zur Nutzung von Hilfsmitteln aus dem Bereich der Künstlichen Intelligenz}
Ich erkläre (Zutreffendes bitte ankreuzen), dass
\begin{enumerate}[label=(\arabic*)]
\item $\square$ ich bei der Erstellung dieser Arbeit keine Hilfsmittel aus dem Bereich der Künstlichen Intelligenz genutzt habe,
\item $\boxtimes$
\begin{itemize}
\item Teile dieser Arbeit unter Verwendung von Hilfsmitteln aus dem Bereich der Künstlichen Intelligenz erstellt wurden (z.\,B. Abbildungen, Text und/oder Quellcode),
\item ich alle Inhalte, die unter Verwendung von Hilfsmitteln aus dem Bereich der Künstlichen Intelligenz erstellt wurden, im Text sowie den zugehörigen Artefakten (z.\,B. Quellcode) an den betreffenden Stellen entsprechend gekennzeichnet habe,
\item mir bewusst bin, dass ich als Autor*in dieser Arbeit die Verantwortung für die in ihr gemachten Angaben und Aussagen trage.
\end{itemize}
\end{enumerate}
\vspace{1cm}

\noindent\parbox[t]{.47\textwidth}{\hrulefill\\Ort, Datum}\hfill
\parbox[t]{.47\textwidth}{\hrulefill\\Unterschrift}
\clearpage
\chapter*{Verwendete KI-Hilfsmittel}
Nach Angabe des Verfassers wurden ChatGPT und teilweise Claude über die unten genannten
Arbeitsschritte hinweg eingesetzt. Der Verfasser hat die Inhalte bearbeitet und trägt die
Verantwortung für ihre Auswahl, Prüfung und Verwendung. Die Angabe beansprucht keine
vollständig ohne KI erzeugte Text-, Code- oder Bildfassung. Die Fußzeilen kennzeichnen
den umfassenden Einsatz an den betreffenden Textstellen; Abbildungsnachweise und
Kommentare in den beigefügten Codeartefakten ergänzen diese Kennzeichnung.

\small
\begin{longtable}{@{}p{.05\textwidth}p{.25\textwidth}p{.64\textwidth}@{}}
\toprule Nr. & Arbeitsschritt & Produkt und Vorgehensweise\\\midrule
1 & $\boxtimes$ Ideen und Konzeption & ChatGPT, teilweise Claude: Unterstützung bei Ideen, Untersuchungsaufbau und Konzeption.\\
2 & $\boxtimes$ Methoden und Modelle & ChatGPT, teilweise Claude: Unterstützung bei Auswahl, Einordnung und Umsetzung.\\
3 & $\boxtimes$ Datensammlung und Analyse & ChatGPT, teilweise Claude: Unterstützung bei Datenaufbereitung und Auswertung.\\
4 & $\boxtimes$ Quellcode & ChatGPT, teilweise Claude: Generierung und Überarbeitung von Code und Notebooks.\\
5 & $\boxtimes$ Textstruktur & ChatGPT, teilweise Claude: Unterstützung bei Gliederung und Aufbau.\\
6 & $\boxtimes$ Textformulierung & ChatGPT, teilweise Claude: Entwürfe und Umformulierungen; Überarbeitung durch den Verfasser.\\
7 & $\boxtimes$ Redaktion & ChatGPT, teilweise Claude; zusätzlich Codex: sprachliche und fachliche Revision.\\
8 & $\boxtimes$ Visualisierungen & ChatGPT, teilweise Claude: Unterstützung bei Abbildungen; Codex: Reproduktion empirischer Diagramme aus Ergebnistabellen.\\
9 & $\boxtimes$ Interpretation und Validierung & ChatGPT, teilweise Claude; Codex: Code-, Quellen- und Zahlenprüfung sowie ergänzende Reproduktionen.\\
10 & $\boxtimes$ Sonstiges & Codex: Rekonstruktion des LaTeX-Projekts, Kompilierung und Layoutprüfung.\\\bottomrule
\end{longtable}
\normalsize
Produkte: ChatGPT -- \url{https://chatgpt.com/}; Claude -- \url{https://claude.ai/};
Codex -- \url{https://openai.com/codex/}. Einzelne Modellversionen früherer Nutzungen
wurden nicht dokumentiert. Der abschließende Codex-Review erfolgte im September 2026.

\vspace{.8cm}
\noindent\parbox[t]{.47\textwidth}{\hrulefill\\Ort, Datum}\hfill
\parbox[t]{.47\textwidth}{\hrulefill\\Unterschrift}
